This section summarizes the results of other research related to your work. In addition to the textual description, please also use tables for this purpose.

\textbf{Note:} The identified research results should be discussed in the discussion section with the results of the thesis.

The Chair of Information Systems and Artificial Intelligence in Marketing uses the \textbf{7\textsuperscript{th} edition of the citation style of the American Psychological Association (APA)}. 

In order to give you an insight into how to cite appropriately in your thesis, the following lines serve as an orientation guideline for the citation system. Adhering to the subse-quent citation rules is vital to crafting a successful and adequate paper. For clarification purposes, several examples are shown in the further procedure.

\subsection{Indirect Citations at the end of a sentence}
If you write a statement or argument in your thesis which is based on the idea or result from another author, you will have to state the author at the end of the sentence. The style of a reference to other scientific work is often not dependent on the type of publication but rather determined by the number of authors.

\subsubsection{Single author}
Theory of mind enables the prediction of others’ behavior beyond familiar patterns, especially in cases involving unfamiliar or mistaken actions \autocite{liszkowski_using_2013}.

\subsubsection{Two authors}
Support for the universality thesis requires that facial expressions of basic emotions are consistently recognized across cultures with high accuracy rates \autocite{nelson_universality_2013}.

\subsubsection{More than two authors}
Young children's early prosocial behavior may stem from their innate interest in others, early imitative tendencies, or emerging capacity to understand others' emotional states \autocite{svetlova_toddlers_2010}.


\subsection{(Indirect) Inline citation style}
If you highlight a result or an idea of some other authors, you can also use inline citations. Make sure that the reference is positioned close to the actual idea of the author(s).

\subsubsection{Single author}
\textcite{izard_many_2010} emphasized that although emotion research has flourished in recent decades, there is still no commonly accepted definition of emotion.

\subsubsection{Two authors}
\textcite{telle_positive_2016} argued that positive empathy also contributes to prosocial behavior, a topic that has received relatively little research attention.

\subsubsection{More than two authors}
\textcite{brownell_peers_2002} found that from around age two, children begin to engage in cooperative play that relies on their developing ability to understand and respond to the mental states of peers.

\subsection{Direct citations}
In addition, for direct citations for which you cite a complete thought you have to indicate this by using quotation marks. For those cases, you also have to specify page numbers. These rules must be applied no matter what you cite (e.g., journals, conferences, books). The style for direct citations is demonstrated by the following examples:

\subsubsection{Direct citation of an entire sentence}
\enquote{[...] children during the second and third years of life do participate in the work of adults.} \autocite[p.~114]{rheingold_little_1982}.

\textbf{Note:} With \enquote{[…]} you highlight that you left out information originally contained in the citation but less relevant for your thesis. The meaning must not be changed by using this technique.

\subsubsection{Direct citation of single information in a sentence}
Already in the earliest stages of development, infants engage in structured interactions with others, often organized as \enquote{coordinated, turn-taking sequences} \autocite[p.~640]{warneken_cooperative_2006}.
